% ===== PRÉAMBULE CANONIQUE — à coller VERBATIM en tête de CHAQUE .tex =====
\documentclass[11pt,a4paper]{article}
\usepackage[utf8]{inputenc}\usepackage[T1]{fontenc}\usepackage{lmodern}
\usepackage[french]{babel}
\usepackage{amsmath,amssymb,mathtools}\usepackage{icomma}
\usepackage[version=4]{mhchem}          % \ce{...} : équations & formules
\usepackage{siunitx}\sisetup{locale=FR} % \qty{}{}, \si{}, \ang{}
\usepackage{chemfig}                    % \chemfig{...} : structures développées
\usepackage{xcolor}\usepackage{colortbl}
\usepackage{tikz}\usepackage{pgfplots}\pgfplotsset{compat=1.18}
\usepackage[most]{tcolorbox}\usepackage{enumitem}\usepackage{array,booktabs}
\usepackage[a4paper,margin=2.2cm]{geometry}
\usetikzlibrary{arrows.meta,calc,angles,quotes,positioning,patterns,backgrounds,shapes.geometric}
\definecolor{primary}{RGB}{222,49,99}\definecolor{primarydark}{RGB}{150,22,55}
\definecolor{defcol}{RGB}{0,114,178}\definecolor{methcol}{RGB}{52,92,148}
\definecolor{excol}{RGB}{0,135,90}\definecolor{attncol}{RGB}{200,40,55}
\tcbset{boxbase/.style={enhanced,breakable,boxrule=0.9pt,arc=2.5pt,
  left=3mm,right=3mm,top=2.3mm,bottom=2.3mm,fonttitle=\bfseries,coltitle=white}}
% --- boîtes TUTORIEL (noms EXACTS, ne pas renommer) ---
\newtcolorbox{cle}[1][]{boxbase,colback=defcol!6,colframe=defcol,
  title={\textsc{À retenir}\ifx\relax#1\relax\relax\else\textmd{ : \itshape #1}\fi}}
\newtcolorbox{methode}[1][]{boxbase,colback=methcol!7,colframe=methcol,sharp corners,
  title={\textsc{Méthode}\ifx\relax#1\relax\relax\else\textmd{ --- \itshape #1}\fi}}
\newtcolorbox{exemplebox}[1][]{boxbase,colback=excol!5,colframe=excol,
  title={\textsc{Exemple}\ifx\relax#1\relax\relax\else\textmd{ : \itshape #1}\fi}}
\newtcolorbox{piege}[1][]{boxbase,colback=attncol!6,colframe=attncol,
  title={\textsc{Piège}\ifx\relax#1\relax\relax\else\textmd{ : \itshape #1}\fi}}
\newtcolorbox{remarque}[1][]{boxbase,colback=black!4,colframe=black!45,
  title={\textsc{Remarque}\ifx\relax#1\relax\relax\else\textmd{ : \itshape #1}\fi}}
% --- boîtes EXERCICE / CORRIGÉ (noms EXACTS, auto-comptées) ---
\newtcolorbox[auto counter]{exo}[1][]{boxbase,colback=primary!4,colframe=primary,
  title={\textsc{Exercice}~\thetcbcounter\ifx\relax#1\relax\relax\else\textmd{ --- \itshape #1}\fi}}
\newtcolorbox[auto counter]{corrige}[1][]{boxbase,colback=excol!4,colframe=excol,
  title={\textsc{Corrigé}~\thetcbcounter\ifx\relax#1\relax\relax\else\textmd{ --- \itshape #1}\fi}}
\newcommand{\R}{\mathbb{R}}\newcommand{\N}{\mathbb{N}}\newcommand{\Z}{\mathbb{Z}}\newcommand{\ds}{\displaystyle}
% (un \usepackage{fancyhdr}+en-tête est possible ; il est ignoré côté web)
% =========================================================================
\begin{document}

\begin{corrige}[nommer un sel neutre ternaire]
\begin{enumerate}[label=\alph*)]
  \item \ce{K2SO4} : sulfate de potassium.
  \item \ce{NaClO} : hypochlorite de sodium.
  \item \ce{Ca(NO3)2} : nitrate de calcium.
\end{enumerate}
\end{corrige}

\begin{corrige}[écrire un sel neutre ternaire]
\begin{enumerate}[label=\alph*)]
  \item carbonate de sodium : \ce{Na+} + \ce{CO3^2-} $\Rightarrow$ \ce{Na2CO3}.
  \item sulfate d'aluminium : \ce{Al^3+} + \ce{SO4^2-} $\Rightarrow$ \ce{Al2(SO4)3}.
  \item permanganate de calcium : \ce{Ca^2+} + \ce{MnO4-} $\Rightarrow$ \ce{Ca(MnO4)2}.
\end{enumerate}
\end{corrige}

\begin{corrige}[quand le métal a une valence variable]
\begin{enumerate}[label=\alph*)]
  \item \ce{Fe(OH)2} : deux \ce{OH-} $\Rightarrow$ fer $+\mathrm{II}$ $\Rightarrow$ hydroxyde de fer(II).
  \item \ce{Fe(OH)3} : trois \ce{OH-} $\Rightarrow$ fer $+\mathrm{III}$ $\Rightarrow$ hydroxyde de fer(III).
  \item nitrate de fer(III) : \ce{Fe^3+} + \ce{NO3-} $\Rightarrow$ \ce{Fe(NO3)3}.
\end{enumerate}
\end{corrige}

\begin{corrige}[nommer un sel acide]
Le \ce{H} restant devient le préfixe « hydrogéno ».
\begin{enumerate}[label=\alph*)]
  \item \ce{NaHCO3} : hydrogénocarbonate de sodium.
  \item \ce{Na2HPO4} : monohydrogénophosphate de sodium.
  \item \ce{NaHSO3} : hydrogénosulfite de sodium.
\end{enumerate}
\end{corrige}

\begin{corrige}[écrire un sel acide]
On croise la valence de l'anion (qui garde son \ce{H}) avec celle du cation.
\begin{enumerate}[label=\alph*)]
  \item hydrogénocarbonate de calcium : \ce{Ca^2+} + \ce{HCO3-} $\Rightarrow$ \ce{Ca(HCO3)2}.
  \item dihydrogénophosphate de potassium : \ce{K+} + \ce{H2PO4-} $\Rightarrow$ \ce{KH2PO4}.
  \item hydrogénosulfate d'ammonium : \ce{NH4+} + \ce{HSO4-} $\Rightarrow$ \ce{NH4HSO4}.
\end{enumerate}
\end{corrige}

\end{document}
